\begin{abstract}

Nesse trabalho buscamos demonstrar experimentalmente a quantização dos níveis de energia de um átomo.
Fizemos isso ao reproduzir o experimento de Franck e Hertz, publicado em 1914 com intuito de testar a 
então recente teoria atômica de N. Bohr.
Nesse experimento, um cátodo aquecido ejeta elétrons dentro de um tubo contendo algum gás.
Os elétrons são acelerados em direção a uma grade que possui potencial positivo.
Alguns deles atravessam a grade e vão em direção a uma placa coletora.
No entanto, a placa possui um pequeno potencial negativo, que repele os elétrons, de modo que apenas 
aqueles com energia suficiente chegam nela.
Nesse trabalho, observamos como a corrente $I$ que chega na placa é afetada ao variar o potencial de excitação $V$, 
o potencial de retardo $V_r$, e a temperatura $T$ no cátodo. 
Encontramos que, se $T$ for suficientemente alto, para valores regularmente espaçados de $V$ as curvas $I(V)$ possuem máximos locais, fato que pode ser explicado pela quantização dos níveis de energia dos átomos do gás.

\end{abstract}
